\documentclass[10pt,a4paper]{amsart}
\usepackage[margin=19mm]{geometry}
\usepackage{amsmath,amssymb,mathtools}
\usepackage[scr=rsfs,cal=euler]{mathalfa}
\usepackage{xfrac}
\usepackage[unicode,hidelinks,bookmarks=false]{hyperref}
\newcommand{\ie}{i.e.\ }
\newcommand{\cf}{cf.\ }
\newcommand{\resp}{respectively\ }
\newcommand{\Subsection}{\subsection}
\providecommand{\nobreakdash}{\nobreak\hbox{-}}
\setlength{\emergencystretch}{2em}
\multlinegap0pt
\usepackage{bm}
\usepackage{bbm}
\usepackage{dsfont}
\usepackage{xspace}
\usepackage{cancel}
\usepackage{mathtools}
\usepackage{upgreek}
\usepackage{amsthm}
\makeatletter
\@ifundefined{lemma}{\newtheorem{lemma}{Lemma}}{}
\makeatother
\def \R{{\mathbb{R}}}
\newcommand{\myparallel}{{\mkern3mu\vphantom{\perp}\vrule depth 0pt\mkern2mu\vrule depth 0pt\mkern3mu}}
\title[Singular limits of anisotropic weak solutions to compressibl]{Singular limits of anisotropic weak solutions to compressible magnetohydrodynamics}
\author{Nicolas Besse, Christophe Cheverry}
\date{Source: arXiv:2506.19784v2 (CC BY 4.0)}
\begin{document}
\maketitle

\noindent\emph{Source and licence.} Singular limits of anisotropic weak solutions to compressible magnetohydrodynamics. Version arXiv:2506.19784v2 (CC BY 4.0), distributed under the Creative Commons Attribution 4.0 International licence, as recorded in the arXiv licence metadata for this exact version and shown on the abstract page https://arxiv.org/abs/2506.19784v2. Full rights notice: licenses/arxiv-2506.19784-CC-BY-4.0.txt.\par\smallskip

\noindent\textit{Editorial scope.} Counted unit: the Lemma whose source label is \texttt{lem:cv-rho-R} in the article (source0.tex lines 1302--1317, source label \texttt{lem:cv-rho-R}), delivered together with the article's own complete original proof of it (source0.tex lines 1319--1357, one \texttt{proof} environment of 39 lines). No mathematics is written, repaired or re-derived: statement and proof are the article's own text line for line. Required context retained uncounted: nrj-ineq-mhdeps (lines 522--525); dissip-D (lines 535--542); ineqC0-R (lines 617--622); defPI-R (lines 641--645); lem:convex (lines 1717--1735); lem:ORC (lines 1743--1755). Every definition, display and label of those lines is retained unchanged. Omitted from this package and not redistributed: 1--521, 526--534, 543--616, 623--640, 646--1301, 1318--1318, 1358--1716, 1736--1742, 1756--2193. Nothing inside a retained excerpt is omitted or altered: every retained body is delivered exactly as the article's lines are written, so no omission receipt of any kind accompanies this package. Citations: the retained excerpts carry no citation command at all, so no bibliography entry of the article is transcribed and no reference is redistributed. Reference adaptation: none was needed and none was made; every reference inside the delivered unit resolves inside it, so no unresolved reference or citation is produced. Printed slips kept literal: no printed word, symbol, hypothesis or number of the counted statement or of its proof was altered. Layout and class: the article's own class is \texttt{amsart} with packages including amsmath, amssymb, amsthm, amsfonts, graphicx, mathrsfs, bbm, xcolor, stmaryrd, upgreek, esint, appendix, mathtools, ulem; the wrapper is the lane's pinned \texttt{amsart} editorial wrapper at 10pt on A4, which restates the theorem-like environments the retained text needs and adds the article's own macro definitions that the retained text needs (\texttt{\textbackslash R}, \texttt{\textbackslash myparallel}), with extra wrapper packages (bm, bbm, dsfont, xspace, cancel, mathtools, upgreek). The delivered numbering is the unit's own. Assets: no retained span contains a figure environment, a graphics-inclusion command, a PDF-page inclusion, a table or a TikZ picture; no image file is copied into this package, the package declares no asset dependency and no image file is redistributed with it. Licence: version arXiv:2506.19784v2 of this article is distributed under the Creative Commons Attribution 4.0 International licence, as recorded in the arXiv licence metadata for this exact version and shown on the abstract page https://arxiv.org/abs/2506.19784v2. A full rights notice is delivered at licenses/arxiv-2506.19784-CC-BY-4.0.txt. Characters: every retained excerpt is pure ASCII, so the delivered engine, XeLaTeX with its default Latin Modern text font, reports no missing character.
\begin{equation}
\label{nrj-ineq-mhdeps}
\mathbbm{E}^\varepsilon(t) + \int_0^t ds \, \mathbbm{D}^\varepsilon(s) \leq  \mathbbm{E}_0^\varepsilon\,, \ \mbox{ a.e. } t \in [0,+\infty)\,,
\end{equation}

\begin{equation}
\label{dissip-D}
 \mathbbm{D}^\varepsilon=
\int_{\Omega} dx\, \big( \mu_\perp^\varepsilon |\nabla_\perp v^\varepsilon|^2 + \mu_\myparallel^\varepsilon |\partial_\myparallel v^\varepsilon|^2 +
\lambda^\varepsilon  |\nabla_\varepsilon \cdot v^\varepsilon|^2 +
\eta_\perp^\varepsilon |\nabla_\perp B^\varepsilon|^2 + \eta_\myparallel^\varepsilon |\partial_\myparallel B^\varepsilon|^2
\big)\,, 
\end{equation}

\begin{equation}
  \label{ineqC0-R}
  \frac 12 \int_{\R^3} dx\, \big( \rho_0^\varepsilon |v_0^\varepsilon|^2 + |B_0^\varepsilon|^2\big) + \frac{a}{\varepsilon^2(\gamma-1)}
  \int_{\R^3} dx\, \big( (\rho_0^\varepsilon)^\gamma - \gamma  \rho_0^\varepsilon 
  + \gamma -1\big) \leq C_0\,.
\end{equation}

\begin{equation}
\label{defPI-R}
\Pi_3(\rho^\varepsilon)=\frac{a}{\varepsilon^2(\gamma-1)}\big( (\rho^\varepsilon)^\gamma
- \gamma \rho^\varepsilon + \gamma - 1\big)\,.
\end{equation}

\begin{lemma}[Convexity properties of the power function]
  \label{lem:convex}
  Let $\bar{x}>0$ and $\gamma >1$ be fixed positive real numbers. For all $ R \in ]\bar{x}, + \infty [ $, there exist positive constants $\nu_i$, $i=1,2,3$,
    depending on $\gamma$, $R$ and $\bar{x}$, such that
  \begin{equation*}
    \begin{aligned}
      x^\gamma  - \gamma x\bar{x}^{\gamma -1} &+ (\gamma-1)\bar{x}^\gamma \\
      & = x^\gamma -\bar{x}^\gamma - \gamma \bar{x}^{\gamma -1}(x-\bar{x}) \geq\\
     & 
    \end{aligned}  
    \left \{
    \begin{aligned}
      \nu_1 \, |x-\bar{x}|^2, & \ \ \ 0 \leq x , \ \  \gamma \geq 2\,, \\
      \nu_2 \, |x-\bar{x}|^2, & \ \ \  0\leq x \leq R,  \ \  1<\gamma < 2\,, \\
      \nu_3 \, |x-\bar{x}|^\gamma, & \ \  \ R < x , \ \   1<\gamma < 2\,.
    \end{aligned}  
    \right.
  \end{equation*}  
\end{lemma}

\begin{lemma}[A criterion for  belonging to the Orlicz spaces $L_2^\gamma(\Omega)$]
  \label{lem:ORC}
  Let $\gamma >1$, $\bar f > 0$ and $\delta >0$ be fixed positive real numbers. Let $ f\in L_{\rm loc}^\gamma(\R^3;\R_+)$ be given. Define
  \[\begin{array}{ll}
    \Uppi_{ \bar f, \gamma } :  L_{\rm loc}^\gamma(\R^3;\R_+) \rightarrow L_{\rm loc}^1(\R^3;\R_+) \,, \quad & \Uppi_{\bar f,\gamma}(f) := f^\gamma - \gamma \bar{f}^{\gamma -1}f +(\gamma -1) {\bar f}^{\gamma} \,, \smallskip \\
    \mathfrak Z_{2,\delta}^{\gamma, \bar f}:   L_{\rm loc}^\gamma (\R^3;\R_+) \rightarrow L_{\rm loc}^1(\R^3;\R_+) \,, \quad & \mathfrak Z_{2,\delta}^{\gamma, \bar f}(f) := |f-\bar f|^2\mathbbm{1}_{\{|f-\bar f|\leq \delta \}} + |f-\bar f|^\gamma\mathbbm{1}_{\{|f-\bar f|> \delta \}} \, .
    \end{array} \]
There exist two constants $\kappa_1$ and $\kappa_2$, depending on $(\bar f, \gamma, \delta)$, such that
    \[
    \kappa_1\, \mathfrak Z_{2,\delta}^{\gamma,\bar f}(f) \leq  \Uppi_{\bar f,\gamma}(f) \leq  \kappa_2\, \mathfrak Z_{2,\delta}^{\gamma, \bar f}(f)\,. 
    \]
  It follows that $\Uppi_{\bar f,\gamma}(f) \in L^1(\R^2)$ if and only if $(f-\bar f) \in L_2^\gamma(\R^3)$.
\end{lemma}

\begin{lemma}
  \label{lem:cv-rho-R}
  There exists a generic constant $C>0$, which may depend on $C_0$,  $a$, and $\gamma$ such that the sequences ${\rho^\varepsilon}$ and  $\varrho^\varepsilon:=(\rho^\varepsilon -1)/\varepsilon$ satisfy the following properties.
  \begin{align*}
    &\| \rho^\varepsilon \|_{ L_{\rm loc}^\infty(\R_+; L_{\rm loc }^\gamma(\R^3))} \leq C\,,
    \quad {\rm and} \quad (\rho^\varepsilon -1) \in L_{\rm loc}^\infty(\R_+; L_2^\gamma \cap H^{-\upalpha}(\R^3))\,, \quad \forall \,\gamma>1 \,, \quad \upalpha \geq 1/2\,,\\
    &\| \rho^\varepsilon -1\|_{ L_{\rm loc}^\infty(\R_+; L^\gamma(\R^3))} \leq  C \varepsilon^{2/\gamma}\,,
    \quad {\rm and} \quad  \| \rho^\varepsilon -1\|_{ L_{\rm loc}^\infty(\R_+; L^2(\R^3))} \leq  C \varepsilon\,,
     \quad \forall\, \gamma \geq 2\,, \\
    & \| \rho^\varepsilon -1\|_{ L_{\rm loc}^\infty(\R_+; L_{\rm loc}^\gamma(\R^3))} \leq  C \varepsilon\,, \quad {\rm for} \ {\rm all } \ \  1<\gamma <2\,,\\
    & \rho^\varepsilon \xrightarrow{\quad} 1 \ \ \mbox{ \rm in } \  L_{\rm loc}^\infty(\R_+; L_2^\gamma \cap  L_{\rm loc}^\gamma(\R^3))\!-\!\mbox{\rm strong}\,, \quad  \forall\, \gamma >1\,,\\
    &  \| \varrho^\varepsilon \|_{ L_{\rm loc}^\infty(\R_+; L_2^\kappa\cap L_{\rm loc }^\kappa\cap H^{-\upalpha}(\R^3))} \leq C\,, 
    \quad \kappa=\min\{2,\gamma\} \,, \quad \upalpha\geq 1/2\\ 
    & \varrho^\varepsilon \xrightharpoonup{\quad} \varrho\ \  \mbox{ \rm in } \  L_{\rm loc}^\infty(\R_+; L_{\rm loc }^\kappa \cap H^{-\upalpha}(\R^3))\!-\!\mbox{\rm weak--}\!\ast\,,  \quad \kappa=\min\{2,\gamma\} \,, \quad \upalpha\geq 1/2\,.
  \end{align*}
\end{lemma}

\begin{proof}
  We start with the first bound of the first assertion of Lemma~\ref{lem:cv-rho-R}. We first claim that $\rho^\varepsilon \in L_{\rm loc}^\infty(\R_+; L_{\rm loc}^\gamma(\R^3))$ if  $\rho^\varepsilon \in L_{\rm loc}^\infty(\R_+; L_{\rm loc}^1(\R^3))$. Indeed, using the convexity of the power function $\R^+\ni x\mapsto x^\gamma$ (with $\gamma >1$), and energy inequality \eqref{nrj-ineq-mhdeps}-\eqref{dissip-D} with the pressure term $\Pi_3$, there exists a constant $ c_0$, such that for any compact set $K\subset \R^3$, $0\leq \int_{K}dx \,\{ (\rho^\varepsilon)^\gamma - \gamma \rho^\varepsilon +\gamma -1 \}\leq  c_0$. Then, $ \int_{K}dx \,(\rho^\varepsilon)^\gamma \leq  c_0 + (\gamma-1)|K|+\gamma \int_{K}dx \,\rho^\varepsilon <\infty$, if  $\rho^\varepsilon \in L_{\rm loc}^\infty(\R_+; L_{\rm loc}^1(\R^3))$. It remains to prove this last property. For this, we consider a test function $\varphi \in \mathscr{C}_c^\infty(\R^3)$, such that $\varphi \geq 0$, and $\varphi \equiv 1$ on a subset $ K $ of its support $ S $. From the mass conservation law and the energy estimate, we obtain
\begin{align*}
  \frac{d}{dt} \int_{S} dx \, \rho^\varepsilon  \varphi  =
  \int_S dx\, \rho^\varepsilon v^\varepsilon \cdot \nabla \varphi 
   \leq \| 2 \nabla \sqrt{\varphi}\|_{L^\infty(\R^3)}\bigg(\int_{\R^3}dx\, \rho^\varepsilon |v^\varepsilon|^2\bigg)
  \bigg(\int_{S}dx\, \rho^\varepsilon \varphi \bigg) \leq \mathfrak c_0
  \bigg(\int_{S} dx\,\rho^\varepsilon \varphi\bigg) \,,
\end{align*}
where the constant $\mathfrak c_0$ depends on $C_0$ and $\varphi$. This inequality leads to $ \int_{K} dx\, |\rho^\varepsilon| \leq \int_{S} dx\, \rho^\varepsilon \varphi \leq \exp(\mathfrak c_0 t)\int_{S} dx\, \rho_0^\varepsilon \varphi $, which shows that $\rho^\varepsilon \in L_{\rm loc}^\infty(\R_+; L_{\rm loc}^1(\R^3))$. It can also be shown from the second and third statements of  Lemma~\ref{lem:cv-rho-R}. Now, from the uniform bound \eqref{ineqC0-R} and energy inequality \eqref{nrj-ineq-mhdeps}-\eqref{dissip-D} with the pressure term $\Pi_3$ defined by \eqref{defPI-R}, we obtain $\Pi_3(\rho^\varepsilon) \in L_{\rm loc}^\infty(\R_+;L^1(\R^3))$ uniformly with respect to $\varepsilon$. From this bound and Lemma~\ref{lem:ORC} (with $f=\rho^\varepsilon$ and $\bar f=1$), we obtain, for any $T\in (0,+\infty)$,
\begin{align}
  \label{eqn:bdze}
 &\sup_{t\in[0,T]}  \int_{\R^3}dx \,\{ |\rho^\varepsilon-1|^2
\mathbbm{1}_{\{|\rho^\varepsilon-1|\leq \delta\}} + |\rho^\varepsilon-1|^\gamma
\mathbbm{1}_{\{|\rho^\varepsilon-1|>\delta\}} \} 
= \sup_{t \in [0,T]}\int_{\R^3} dx\, \mathfrak Z_{2,\delta}^{\gamma,1}(\rho^\varepsilon) \nonumber \\
& \hspace{2.5cm} \leq \sup_{t \in [0,T]} \frac{1}{\kappa_1}\int_{\R^3} dx\, \Uppi_{1,\gamma}(\rho^\varepsilon)
= \sup_{t \in [0,T]} \frac{(\gamma-1)\varepsilon^2}{\kappa_1  a} \int_{\R^3} dx\, \Pi_3(\rho^\varepsilon)
\leq\frac{ C_0(\gamma-1)\varepsilon^2}{\kappa_1  a}\,.
\end{align}
This inequality implies the bound in $L_{\rm loc}^\infty(\R_+; L_2^\gamma(\R^3))$ in the second part of the first  statement of Lemma~\ref{lem:cv-rho-R}, and also the strong convergence of $\rho^\varepsilon$ to $1$ in $L_{\rm loc}^\infty(\R_+, L_2^\gamma(\R^3))$ in the fourth statement of Lemma~\ref{lem:cv-rho-R}. To complete the proof of the second part of the first  statement of Lemma~\ref{lem:cv-rho-R}, we have to show that $L_2^\gamma(\R^3) \hookrightarrow H^{-\upalpha}(\R^3)$, with $\upalpha \geq 1/2$. This embedding is obvious for $\gamma=2$. For $\gamma \neq 2$, from the definition of the space  $L_2^\gamma(\R^3)$, the density $\rho^\varepsilon(t)-1$ can be splitted into parts $d_1(t)^\varepsilon:=(\rho^\varepsilon(t) -1) \mathbbm{1}_{\{|{\rho^\varepsilon(t)-1}|\leq \delta\}} \in L^2(\R^3)$  and $d_2^\varepsilon(t):=(\rho^\varepsilon(t) -1) \mathbbm{1}_{\{|{\rho^\varepsilon(t)-1}|> \delta\}} \in L^\gamma(\R^3)$. The part $d_1^\varepsilon(t)$ is obviously in $H^{-\upalpha}(\R^3)$ with $\upalpha \geq 1/2$.  For the part $d_2^\varepsilon(t)$ we proceed as follows. By Sobolev embeddings we have $H^\upalpha(\R^3)  \hookrightarrow L^{\gamma'}(\R^3)$ with $1/\gamma +1/\gamma'=1$, and $\upalpha>1/2\, $ since $\gamma>3/2$. By duality we then obtain $d_2^\varepsilon(t) \in L^\gamma(\R^3)  \hookrightarrow H^{-\upalpha}(\R^3)$. Therefore $(\rho^\varepsilon(t) -1) \in  H^{-\upalpha}(\R^3)$, with $\upalpha \geq 1/2$. We continue with the second assertion of  Lemma~\ref{lem:cv-rho-R}. Estimate \eqref{eqn:bdze} and the inequality $|\rho^\varepsilon -1|^2 \geq |\rho^\varepsilon-1|^\gamma$, for $\gamma\geq 2$ and $|\rho^\varepsilon-1| \leq \delta <1$, imply the first part of the second statement of Lemma~\ref{lem:cv-rho-R}. The second part of this second statement appeals to Lemma~\ref{lem:convex}. Indeed, from this lemma, for $\gamma \geq 2$, we obtain $|\rho^\varepsilon -1|^2 \leq \varepsilon^2(\gamma-1)/(\nu_1  a)\Pi_3(\rho^\varepsilon)$, which gives $\| \rho^\varepsilon -1\|_{ L_{\rm loc}^\infty(\R_+; L^2(\R^3))} \leq  \varepsilon \sqrt{C_0 (\gamma-1)/(\nu_1  a)}$. We continue with the third assertion of Lemma~\ref{lem:cv-rho-R}. Using Cauchy--Schwarz inequality and estimate \eqref{eqn:bdze}, we obtain, for any compact set $K\subset \R^3$, and any $T\in (0,+\infty)$,
\begin{align*}
 \sup_{t\in [0,T]} \|\rho^\varepsilon(t) -1\|_{L^\gamma(K)}^\gamma
  & \leq |K|^{1-\gamma/2}  \sup_{t\in [0,T]} \bigg(\int_K dx\, |\rho^\varepsilon -1|^2
  \mathbbm{1}_{\{|{\rho^\varepsilon-1}|\leq \delta\}} \bigg)^{\gamma/2} \\
  & \quad + \sup_{t\in [0,T]} \int_K dx\, |\rho^\varepsilon -1|^\gamma
\mathbbm{1}_{\{|{\rho^\varepsilon-1}|> \delta\}} \\
& \leq C(|K|, C_0,  a,  \gamma) (\varepsilon^\gamma + \varepsilon^2) \,,
\end{align*}
which justifies the third assertion of Lemma~\ref{lem:cv-rho-R}. Then, the strong convergence of $\rho^\varepsilon$ to $1$ in $L_{\rm loc}^\infty(\R_+, L_{\rm loc}^\gamma(\R^3))$ in the fourth statement of Lemma~\ref{lem:cv-rho-R}, is obtained from the first part of the second statement of Lemma~\ref{lem:cv-rho-R}, and the third statement of Lemma~\ref{lem:cv-rho-R}. We continue with the fifth statement of Lemma~\ref{lem:cv-rho-R}. The uniform bound in $L_{\rm loc}^\infty(\R_+; L_{\rm loc }^\kappa(\R^3))$ for $\varrho^\varepsilon$ comes from the second part of the second statement of Lemma~\ref{lem:cv-rho-R}, and the third statement of Lemma~\ref{lem:cv-rho-R}. For the uniform  bound of  $\varrho^\varepsilon$ in $L_{\rm loc}^\infty(\R_+; L_2^\kappa(\R^3))$, we distinguish two cases according to the value of $\gamma$. For $\gamma \geq 2$, since $\kappa=2$, the second part of the second statement of Lemma~\ref{lem:cv-rho-R} implies the fifth one. For $1<\gamma <2$, estimate \eqref{eqn:bdze} leads to
\[
\sup_{t\in[0,T]}  \int_{\R^3}dx \, \Big\{ \Big|\frac{\rho^\varepsilon-1}{\varepsilon}\Big|^2
\mathbbm{1}_{\{|\frac{\rho^\varepsilon-1}{\varepsilon}|\leq \frac{ \delta}{\varepsilon}\}} + 
\varepsilon^{\gamma-2}\Big|\frac{\rho^\varepsilon-1}{\varepsilon}\Big|^\gamma
\mathbbm{1}_{\{|\frac{\rho^\varepsilon-1}{\varepsilon}|> \frac{ \delta}{\varepsilon}\}}
\Big\} \leq \frac{C_0(\gamma-1)}{\kappa_1 a}\,,
\]
for any $T\in (0,+\infty)$. This last  estimate and  inequality $\varepsilon^{\gamma-2}>1$, imply the bound of $\varrho^\varepsilon$ in $L_{\rm loc}^\infty(\R_+; L_2^\kappa(\R^3))$. To complete the fifth assertion of Lemma~\ref{lem:cv-rho-R}, we observe, as above, that we have the embedding $L_2^\kappa(\R^3) \hookrightarrow H^{-\upalpha}(\R^3)$, with $\upalpha \geq 1/2$ and $\kappa=\min\{2,\gamma\}$. Finally, the sixth statement of Lemma~\ref{lem:cv-rho-R} is obtained from the fifth assertion and weak compactness properties. This ends the proof of Lemma~\ref{lem:cv-rho-R}.
\end{proof}  

\end{document}
